% FILE: 02_lineage_and_context.tex
% Project: prompts — Downstream Manufacturing Prompt Corpus

\section{Lineage and Context}
\label{sec:prompts-lineage}

\subsection{2016 Language-Model Lesson}
\label{sec:prompts-lineage-2016}

The thesis lineage begins with a 2016 language-model experiment whose finding was
that local text imitation does not conserve consequence. A model that could reproduce
the surface form of a document could not be trusted to reproduce the causal structure
that made the document's claims binding. This lesson is directly instantiated in the
prompt corpus: a directive that says ``implement Inductive Miner'' without specifying
what the return type must be, what fixtures must pass, and what authority authorizes
the work is a surface-form instruction that does not conserve consequence. The
downstream engineering team could implement an algorithm that produces a result with
the right shape but the wrong mathematical guarantees, and no artifact would capture
the failure.

The corpus answers the 2016 lesson by requiring that every directive specify not just
the surface behavior but the structural consequence: the Rust type that encodes the
guarantee (\texttt{TypedProcessTree}), the compile-fail fixture that proves the
guarantee cannot be bypassed, and the authority source that makes the mandate
binding. The directive is not a natural-language suggestion; it is a manufacturing
specification whose acceptance criteria are mechanically verifiable.

\subsection{Enterprise Process Gap}
\label{sec:prompts-lineage-enterprise}

The enterprise process gap is the observation that process knowledge in organizations
exists in two incompatible forms: declarative models (BPMN diagrams, process
documentation, policy statements) and execution evidence (event logs, audit trails,
transaction records). The gap arises because there is no formal mechanism to compare
the two. An organization can claim conformance to a process model without any
evidence that live execution actually follows it, and a software system can pass all
of its unit tests without any proof that the process it executes is sound.

The prompt corpus addresses the enterprise process gap at the engineering layer. The
Blue River Dam lifecycle directive (\texttt{downstream\_blue\_river\_dam\_lifecycle\_authority.md})
requires that every state transition produce a cryptographically signed proof receipt.
Gate 1 (Design) requires a soundness receipt proving the Petri net $N = (P, T, F)$
satisfies all four WF-net axioms. Gate 3 (Monitoring Operations) requires that every
live trace produce a fitness receipt, with explicit admissibility rules:
fitness $\geq 0.95$ admits automatically; $0.85 \leq f < 0.95$ requires an Executive
Board override signature; $f < 0.85$ triggers a \texttt{RefusalReport} and the trace
is rejected. These requirements close the enterprise process gap at the implementation
level by making it impossible for a process to claim ALIVE status without a complete
receipt chain.

\subsection{Relationship to the Chatman Equation}
\label{sec:prompts-lineage-chatman}

The Chatman Equation $\kappa \circ \rho \circ \alpha \circ \mu$ is the compositional
identity of the process intelligence architecture. The four operators are:
$\mu$ (Manufacture) --- initial admission of the WASM payload via
\texttt{wasm4pm-compat}; $\alpha$ (Attest) --- cryptographic validation and typestate
enforcement; $\rho$ (Receipt) --- BLAKE3 receipt generation; and $\kappa$ (Keeper)
--- final execution state within \texttt{wasm4pm}. The equation asserts that a
process execution claim is only valid when all four operators have composed
successfully in sequence.

The \texttt{execution-plans/blue-river-dam-bridge.md} blueprint makes this equation
physically concrete. It specifies a three-phase cryptographic linkage: in the
Admission Phase, \texttt{wasm4pm-compat} generates a preliminary AST and a BLAKE3
hash of the raw binary stream; in the Attestation and Receipt phase, the
\texttt{wasm4pm} engine consumes the preliminary hash and AST, executes a
deterministic compilation pass, and synthesizes a Receipt object containing
\texttt{H\_orig} (BLAKE3 hash of original payload), \texttt{H\_exec} (BLAKE3 hash
of compiled artifact), and \texttt{Sig\_adm} (Ed25519 signature from the compat
admission enclave); in the Execution Phase, the sandbox verifies \texttt{Sig\_adm}
and \texttt{H\_exec} before instantiation. Typestate enforcement via Rust's phantom
type pattern makes the Receipt object the only lawful path to instantiation.

\subsection{Relationship to Receipts and Replay}
\label{sec:prompts-lineage-receipts}

The receipts-and-replay design principle is that every claim about process behavior
must be backed by a cryptographic receipt that can be independently verified and
replayed to confirm the claim. The prompt corpus extends this principle to the
manufacturing directive layer itself. The \texttt{ProcessIntelligenceVerificationReceipt}
JSON schema specified in \texttt{downstream\_ggen\_projection\_integration.md}
requires: a \texttt{receipt\_id} (UUID), \texttt{target\_log\_hash} (SHA-256 of the
source event log), \texttt{process\_model\_hash} (SHA-256 of the model), fitness and
precision values computed by \texttt{wasm4pm}, and a \texttt{validator\_signature}
(base64-encoded Ed25519). A receipt is refused if any required field is missing, if
the timestamp is future-dated or stale beyond 90 days, or if the Ed25519 signature
fails verification against the \texttt{RoleKeyRegistry}.

The ggen integration directive formalizes the replay direction: a board claim is only
admissible if the supporting receipt can be re-verified by replaying the admission
pipeline with the same log hash and model hash and obtaining the same signature. This
makes the M\&A claim manufacturing process fully auditable and replayable, closing
the loop from execution evidence to capital-flow assertion.
